\documentclass[11pt]{article}
\usepackage{amsmath,amssymb,amsthm}
\usepackage[margin=1in]{geometry}
\usepackage{hyperref}
\usepackage{booktabs}
\usepackage{tikz-cd}

\newtheorem{theorem}{Theorem}
\newtheorem{proposition}{Proposition}
\newtheorem{conjecture}{Conjecture}
\newtheorem{definition}{Definition}
\newtheorem{remark}{Remark}

\newcommand{\KL}{\mathrm{KL}}
\newcommand{\OT}{\mathrm{OT}}
\newcommand{\VFE}{\mathrm{VFE}}
\newcommand{\GL}{\mathrm{GL}}
\newcommand{\SO}{\mathrm{SO}}
\newcommand{\tr}{\mathrm{tr}}
\newcommand{\R}{\mathbb{R}}
\newcommand{\E}{\mathbb{E}}

\title{Toward a Unified Theory of Gauge-VFE and Optimal Transport}
\author{Research Notes}
\date{March 2026}

\begin{document}
\maketitle

\begin{abstract}
We outline a roadmap for unifying the gauge-covariant variational free energy (Gauge-VFE) framework with optimal transport (OT) theory. The two frameworks share deep structural parallels---both concern the geometry of probability distributions, both define notions of ``cost'' for moving mass/belief between configurations, and both give rise to natural gradient flows. We identify precise mathematical contact points, open problems, and a phased research program toward a unified theory in which Gauge-VFE attention emerges as a gauge-equivariant optimal transport problem on a statistical fiber bundle.
\end{abstract}

\tableofcontents

%=============================================================================
\section{Executive Summary: Why These Theories Should Unify}
%=============================================================================

The Gauge-VFE framework and optimal transport theory operate on the same mathematical terrain---the geometry of probability measures---but approach it from complementary directions:

Gauge-VFE starts from a \emph{fiber bundle} perspective: agents carry local beliefs $q_i = \mathcal{N}(\mu_i, \Sigma_i)$ in local gauge frames $\phi_i$, and the free energy $F$ penalizes KL divergence between gauge-transported beliefs. Attention weights $\beta_{ij} = \text{softmax}(-\KL(q_i \| \Omega_{ij} q_j)/\kappa)$ emerge as the mechanism for belief alignment.

Optimal transport starts from a \emph{coupling} perspective: given source measure $\mu$ and target measure $\nu$, find the coupling $\pi \in \Pi(\mu,\nu)$ that minimizes the total transport cost $\int c(x,y)\,d\pi(x,y)$. The Wasserstein distance $W_p$ metrizes the space of measures, and gradient flows in Wasserstein space characterize diffusion, aggregation, and information dynamics.

The core claim of this document is that these are two views of the same underlying structure, and a proper unification yields:
\begin{enumerate}
    \item A transport-theoretic interpretation of attention as an optimal coupling,
    \item A gauge-theoretic interpretation of the Monge map as parallel transport,
    \item A unified variational principle that subsumes both frameworks,
    \item Computational benefits from OT algorithms (Sinkhorn, etc.) applied to gauge-VFE dynamics.
\end{enumerate}


%=============================================================================
\section{Background: The Two Frameworks}
%=============================================================================

\subsection{Gauge-VFE (This Project)}

The full variational free energy over $N$ agents on a statistical fiber bundle:
\begin{equation}\label{eq:vfe}
F = \underbrace{\alpha \sum_i \KL(q_i \| p_i)}_{\text{self-coupling}}
  + \underbrace{\lambda_\beta \sum_{i,j} \beta_{ij} \KL(q_i \| \Omega_{ij} q_j)}_{\text{belief alignment}}
  + \underbrace{\text{CE}(\mu, y)}_{\text{observation}}
  + \underbrace{\lambda_p \sum_{i,j} \KL(p_i \| \Omega_{ij} p_j)}_{\text{prior alignment}}
\end{equation}

Key structures:
\begin{itemize}
    \item \textbf{Statistical fiber bundle}: base manifold $C$ (token positions), fiber $\mathcal{P}$ (Gaussian beliefs), structure group $\GL(K)$.
    \item \textbf{Gauge transport}: $\Omega_{ij} = \exp(\phi_i)\exp(-\phi_j) \in \GL^+(K)$, acting on beliefs via pushforward $(\mu, \Sigma) \mapsto (\Omega\mu, \Omega\Sigma\Omega^\top)$.
    \item \textbf{Attention}: $\beta_{ij} = \text{softmax}_j(-\KL(q_i \| \Omega_{ij}q_j)/\kappa)$.
    \item \textbf{GL(K) gauge invariance}: $\KL$ is invariant under the full general linear group---the $(\det\Omega)^2$ factors cancel in the log-determinant ratio.
\end{itemize}

\subsection{Optimal Transport}

Given probability measures $\mu, \nu$ on spaces $\mathcal{X}, \mathcal{Y}$ and cost function $c: \mathcal{X} \times \mathcal{Y} \to \R_+$:

\textbf{Kantorovich formulation} (relaxed):
\begin{equation}\label{eq:kantorovich}
\OT_c(\mu,\nu) = \inf_{\pi \in \Pi(\mu,\nu)} \int c(x,y)\,d\pi(x,y)
\end{equation}
where $\Pi(\mu,\nu)$ is the set of couplings with marginals $\mu$ and $\nu$.

\textbf{Monge formulation} (deterministic):
\begin{equation}
\OT_c(\mu,\nu) = \inf_{T: T_\#\mu = \nu} \int c(x, T(x))\,d\mu(x)
\end{equation}

\textbf{Wasserstein-2 distance}: With $c(x,y) = \|x-y\|^2$,
\begin{equation}
W_2^2(\mu,\nu) = \inf_{\pi \in \Pi(\mu,\nu)} \int \|x - y\|^2\,d\pi(x,y)
\end{equation}

\textbf{Entropic regularization} (Sinkhorn):
\begin{equation}\label{eq:sinkhorn}
\OT_c^\varepsilon(\mu,\nu) = \inf_{\pi \in \Pi(\mu,\nu)} \int c(x,y)\,d\pi(x,y) + \varepsilon \KL(\pi \| \mu \otimes \nu)
\end{equation}

\textbf{Gaussian OT} (Bures-Wasserstein): For $\mu = \mathcal{N}(m_1, S_1)$, $\nu = \mathcal{N}(m_2, S_2)$:
\begin{equation}\label{eq:bures}
W_2^2(\mu,\nu) = \|m_1 - m_2\|^2 + \tr\!\left(S_1 + S_2 - 2(S_1^{1/2} S_2 S_1^{1/2})^{1/2}\right)
\end{equation}
This is the \textbf{Bures metric} on the space of Gaussian distributions.


%=============================================================================
\section{Structural Correspondences}
%=============================================================================

We now identify precise mathematical contact points between the two frameworks.

\subsection{Attention as Soft Optimal Coupling}

The attention matrix $\beta_{ij}$ in Gauge-VFE has the structure of an entropically regularized optimal transport plan.

\begin{proposition}[Attention $\leftrightarrow$ Sinkhorn]\label{prop:attention-sinkhorn}
The Gauge-VFE attention weights
\[
\beta_{ij} = \frac{\exp(-\KL(q_i \| \Omega_{ij}q_j)/\kappa)}{\sum_k \exp(-\KL(q_i \| \Omega_{ik}q_k)/\kappa)}
\]
are the row-normalized solution to an entropy-regularized transport problem with:
\begin{itemize}
    \item Cost matrix $C_{ij} = \KL(q_i \| \Omega_{ij}q_j)$,
    \item Temperature $\varepsilon = \kappa$,
    \item Source marginal: uniform over queries (row normalization via softmax),
    \item Target marginal: unconstrained (single-sided normalization).
\end{itemize}
\end{proposition}

The standard Sinkhorn algorithm enforces \emph{both} marginal constraints via alternating row/column normalization. Standard attention enforces only the row constraint (softmax over keys). This is not an accident---it reflects the asymmetry of KL divergence. A fully doubly-stochastic attention (enforcing both marginals) would correspond to a proper Sinkhorn solution and would impose a conservation law: total attention received equals total attention given.

\begin{remark}[Sinkhorn Attention]
Sinkhorn attention (doubly stochastic normalization of the attention matrix) has been explored in the literature. Within Gauge-VFE, this would correspond to enforcing a measure-preservation constraint: the attention coupling preserves the total ``belief mass'' across the system. This is physically natural if we interpret belief alignment as a transport of information that conserves total information content.
\end{remark}

\subsection{Gauge Transport as Monge Map}

The gauge transport operator $\Omega_{ij}$ pushes belief $q_j$ forward into frame $i$:
\[
(\Omega_{ij})_\# q_j = \mathcal{N}(\Omega_{ij}\mu_j,\; \Omega_{ij}\Sigma_j\Omega_{ij}^\top)
\]

This is precisely a \emph{linear Monge map}: the deterministic transport $T_{ij}(x) = \Omega_{ij}x$ that pushes $q_j$ toward $q_i$'s frame.

\begin{proposition}[Gauge Transport $\leftrightarrow$ Linear Monge Map]\label{prop:gauge-monge}
For Gaussian beliefs $q_i = \mathcal{N}(\mu_i, \Sigma_i)$ and $q_j = \mathcal{N}(\mu_j, \Sigma_j)$, the optimal transport map $T^*: \R^K \to \R^K$ satisfying $T^*_\# q_j = q_i$ (for quadratic cost) is the affine map:
\[
T^*(x) = \mu_i + A(\Sigma_j, \Sigma_i)(x - \mu_j)
\]
where $A(\Sigma_j, \Sigma_i) = \Sigma_j^{-1/2}(\Sigma_j^{1/2}\Sigma_i\Sigma_j^{1/2})^{1/2}\Sigma_j^{-1/2}$ is the Brenier map.

The gauge transport $\Omega_{ij}$ approximates this when:
\begin{enumerate}
    \item Covariances are isotropic: $\Sigma_i = \sigma^2 I$ $\Rightarrow$ $A = I$ and transport reduces to mean translation,
    \item The gauge field is adapted to the covariance structure: $\Omega_{ij} \approx A(\Sigma_j, \Sigma_i)$.
\end{enumerate}

In the general non-isotropic case, $\Omega_{ij}$ is a \emph{constrained} Monge map---it is the best \emph{linear} transport within the $\GL(K)$ gauge orbit, rather than the globally optimal Brenier map.
\end{proposition}

This gives a clear interpretation: the gauge group $\GL(K)$ parameterizes the space of admissible transport maps, and the VFE selects the optimal one within this family.

\subsection{KL Divergence vs.\ Wasserstein Distance}

The two frameworks use different ``costs'' for comparing distributions:

\begin{center}
\begin{tabular}{lcc}
\toprule
& \textbf{Gauge-VFE} & \textbf{Optimal Transport} \\
\midrule
Cost & $\KL(q_i \| \Omega_{ij}q_j)$ & $W_2^2(q_i, q_j)$ \\
Symmetry & Asymmetric & Symmetric \\
Triangle ineq. & No & Yes (metric) \\
Gauge invariance & $\GL(K)$-invariant & Isometry-invariant \\
Gaussian formula & Analytic (trace + log-det) & Bures metric (matrix square root) \\
Gradient flow & Fisher-Rao natural gradient & Wasserstein gradient (continuity eq.) \\
\bottomrule
\end{tabular}
\end{center}

For Gaussians with $\Sigma_i = \sigma^2 I$ (isotropic), both reduce to squared Euclidean distance on means:
\begin{align}
\KL(q_i \| q_j) &= \frac{1}{2\sigma^2}\|\mu_i - \mu_j\|^2 \\
W_2^2(q_i, q_j) &= \|\mu_i - \mu_j\|^2
\end{align}
differing only by the precision factor $1/(2\sigma^2)$. This is the isotropic limit where standard dot-product attention is recovered.

\subsection{The Bures--Fisher Bridge}

A deep connection exists between the Bures-Wasserstein metric (from OT) and the Fisher-Rao metric (used in Gauge-VFE for natural gradients).

For Gaussian distributions, consider infinitesimal perturbations $(\delta\mu, \delta\Sigma)$:

\textbf{Fisher-Rao metric:}
\begin{equation}
ds_{\text{FR}}^2 = \delta\mu^\top \Sigma^{-1} \delta\mu + \frac{1}{2}\tr(\Sigma^{-1}\delta\Sigma\,\Sigma^{-1}\delta\Sigma)
\end{equation}

\textbf{Bures-Wasserstein metric:}
\begin{equation}
ds_{\text{BW}}^2 = \|\delta\mu\|^2 + \frac{1}{4}\tr\!\left(\delta\Sigma\,\Sigma^{-1}\delta\Sigma\,\Sigma^{-1} + \delta\Sigma^2\Sigma^{-2}\right)
\end{equation}

At $\Sigma = \sigma^2 I$, both metrics become proportional. More importantly, both are Riemannian metrics on the Gaussian manifold, and the identity map between $(G_{\text{FR}}, G_{\text{BW}})$ is a bi-Lipschitz equivalence with dimension-dependent constants.

The Gauge-VFE pullback metrics (the ``it from bit'' construction in \texttt{pullback\_metrics.py}) pull back $g_{\text{FR}}$ to the base manifold. An analogous construction pulling back $g_{\text{BW}}$ would yield an emergent geometry measuring Wasserstein-optimal transport distances rather than Fisher-Rao statistical distances on the base.

\begin{conjecture}[Dual Emergent Geometries]
The pullback of the Fisher-Rao metric to the base manifold $C$ gives the ``epistemic geometry'' (information distance). The pullback of the Bures-Wasserstein metric gives the ``transport geometry'' (physical displacement distance). These two geometries are related by the gauge field: the Fisher-Rao geometry is the transport geometry ``dressed'' by the gauge connection.
\end{conjecture}


%=============================================================================
\section{The Unification Program}
%=============================================================================

We propose a phased research program toward a unified Gauge-VFE-OT theory.

\subsection{Phase 1: Replace KL with Bures-Wasserstein in Attention}

The most direct connection: replace the KL-divergence cost in attention with the Bures-Wasserstein distance.

\textbf{Current (KL-based attention):}
\begin{equation}
\beta_{ij} = \text{softmax}_j\!\left(-\KL(q_i \| \Omega_{ij}q_j)/\kappa\right)
\end{equation}

\textbf{Proposed (Bures-Wasserstein attention):}
\begin{equation}
\beta_{ij}^{\text{BW}} = \text{softmax}_j\!\left(-W_2^2(q_i, (\Omega_{ij})_\#q_j)/\kappa\right)
\end{equation}

For Gaussians, $W_2^2$ has closed form (Eq.~\ref{eq:bures}). This replaces the asymmetric KL with a proper metric while retaining gauge transport.

\textbf{Advantages:}
\begin{itemize}
    \item $W_2^2$ is a genuine metric (symmetric, satisfies triangle inequality).
    \item The Bures metric has a natural geometric interpretation as the geodesic distance on the manifold of Gaussians under the Wasserstein geometry.
    \item Gradients of $W_2^2$ are well-understood and relate to optimal transport maps.
\end{itemize}

\textbf{Key question:} Does GL(K) gauge invariance survive? The Bures metric is invariant under \emph{orthogonal} transformations $O \in O(K)$, not the full GL(K). This would break the GL(K) gauge symmetry down to $O(K)$, which is more restrictive (requires orthogonal gauge transport). This may or may not be desirable---it trades computational simplicity (no re-orthogonalization) for metric-space structure.

\textbf{Resolution path:} Define a \emph{gauge-covariant Bures distance}:
\begin{equation}
W_{2,\Omega}^2(q_i, q_j) \coloneqq W_2^2(q_i, (\Omega_{ij})_\#q_j)
\end{equation}
This is gauge-invariant by construction: the transport $\Omega_{ij}$ absorbs the gauge transformation before the Bures comparison. The residual symmetry is $O(K) \cap \text{Stab}(\Omega_{ij})$.

\subsection{Phase 2: Sinkhorn Attention (Doubly Stochastic)}

Enforce both marginal constraints on the attention matrix to obtain a proper transport plan.

\textbf{Algorithm:}
\begin{enumerate}
    \item Compute cost matrix $C_{ij} = \KL(q_i \| \Omega_{ij}q_j)$ (or $W_2^2$).
    \item Initialize $K = \exp(-C/\varepsilon)$ (Gibbs kernel).
    \item Sinkhorn iterations: alternate row and column normalization until convergence.
    \item Result: $\beta^* = \text{diag}(u)\,K\,\text{diag}(v)$ with prescribed marginals.
\end{enumerate}

This gives attention matrices that satisfy a \emph{conservation law}: each token receives a fixed total attention (e.g., uniform marginals mean every token is equally ``heard''). This is natural from a transport perspective---information is neither created nor destroyed, only moved.

\textbf{Connection to existing work:} Sinkhorn transformers (Tay et al., 2020) apply doubly stochastic normalization but without the gauge-theoretic cost function. The Gauge-VFE version would use the geometrically principled $\KL$-based cost.

\subsection{Phase 3: Wasserstein Gradient Flows for Belief Dynamics}

Currently, beliefs evolve via Fisher-Rao natural gradient descent:
\begin{equation}
\dot{\mu}_i = -\Sigma_i \nabla_{\mu_i} F \qquad \text{(natural gradient)}
\end{equation}

The alternative is a \emph{Wasserstein gradient flow}: the belief distribution $q_i(t)$ evolves according to the continuity equation
\begin{equation}
\partial_t q_i + \nabla \cdot (q_i\, v_i) = 0
\end{equation}
where the velocity field $v_i = -\nabla \frac{\delta F}{\delta q_i}$ is the Wasserstein gradient of the free energy.

For Gaussians, this reduces to coupled ODEs for $(\mu_i, \Sigma_i)$. The key difference from the Fisher-Rao flow:

\begin{center}
\begin{tabular}{lcc}
\toprule
& \textbf{Fisher-Rao (current)} & \textbf{Wasserstein (proposed)} \\
\midrule
Mean update & $\dot{\mu} = -\Sigma\nabla_\mu F$ & $\dot{\mu} = -\nabla_\mu F$ \\
Covariance update & $\dot{\Sigma} = -2\Sigma\,\text{sym}(\nabla_\Sigma F)\,\Sigma$ & Riccati-type ODE \\
Geometry & Reparametrization-invariant & Physically ``moves mass'' \\
Convergence & Fast for peaked distributions & Uniform across scales \\
\bottomrule
\end{tabular}
\end{center}

\begin{remark}
The JKO scheme (Jordan-Kinderlehrer-Otto) interprets diffusion equations as Wasserstein gradient flows of the free energy functional. The VFE belief dynamics could be reformulated as a JKO-type scheme, where each belief step minimizes:
\[
q_i^{(t+1)} = \arg\min_q \left\{ F[q] + \frac{1}{2\tau}W_2^2(q, q_i^{(t)}) \right\}
\]
This proximal operator interpretation connects VFE inference to implicit OT time-stepping.
\end{remark}

\subsection{Phase 4: Gauge-Equivariant Optimal Transport}

The deepest level of unification: formulate the \emph{entire VFE} as an optimal transport problem on the statistical fiber bundle.

\begin{definition}[Gauge-Equivariant Transport Problem]
Let $\mathcal{E} = (C, \mathcal{P}, \GL(K))$ be the statistical fiber bundle. Given belief configurations $\{q_i\}$ and prior configurations $\{p_i\}$, the gauge-equivariant optimal transport problem is:
\begin{equation}\label{eq:gauge-ot}
\OT_{\text{gauge}}(\{q\}, \{p\}) = \inf_{\substack{\pi \in \Pi(q, p) \\ \Omega \in \GL(K)^{N \times N}}} \sum_{i,j} \pi_{ij}\, c_\Omega(q_i, p_j)
\end{equation}
where $c_\Omega(q_i, p_j) = W_2^2(q_i, (\Omega_{ij})_\# p_j)$ and $\Omega_{ij}$ satisfies the cocycle condition $\Omega_{ij}\Omega_{jk} = \Omega_{ik}$.
\end{definition}

This jointly optimizes:
\begin{enumerate}
    \item The \textbf{coupling} $\pi$ (which beliefs align with which priors---the attention pattern),
    \item The \textbf{gauge field} $\Omega$ (how to transport between frames---the connection).
\end{enumerate}

The existing VFE (Eq.~\ref{eq:vfe}) is a relaxation of this problem:
\begin{itemize}
    \item The coupling $\pi_{ij} = \beta_{ij}$ is not jointly optimized but derived from the cost via softmax (entropic regularization).
    \item The gauge field $\Omega_{ij} = \exp(\phi_i)\exp(-\phi_j)$ is parameterized (not freely optimized).
    \item The cost is $\KL$ rather than $W_2^2$.
\end{itemize}

\subsection{Phase 5: Renormalization as Multi-Scale OT}

The RG flow in Gauge-VFE (meta-agent emergence via spectral clustering on $\beta$) has a direct OT interpretation: \emph{hierarchical optimal transport}.

At scale $\zeta = 0$, tokens interact via fine-grained transport $\pi^{(0)}$. Coarse-graining to scale $\zeta = 1$ (meta-agents) corresponds to solving a reduced OT problem on the clustered measures. This is precisely the structure of \textbf{multi-scale optimal transport} (Gerber \& Maggioni, 2017), where one builds a tree of OT problems at successively coarser scales.

\begin{conjecture}[RG Flow $=$ Multi-Scale OT]
The RG fixed point of Gauge-VFE dynamics---where modularity $Q(\beta)$ stabilizes and the effective rank plateaus---coincides with the scale at which the hierarchical OT cost is minimized for the given observation.
\end{conjecture}


%=============================================================================
\section{Concrete Mathematical Results to Prove}
%=============================================================================

\subsection{Theorem Targets}

\begin{enumerate}

\item \textbf{KL--Bures equivalence at isotropic limit.}
\emph{Claim:} For $\Sigma_i = \sigma^2 I$, $\KL(q_i \| \Omega_{ij}q_j) = \frac{1}{2\sigma^2}W_2^2(q_i, (\Omega_{ij})_\#q_j) + \text{const}$, so the two attention mechanisms coincide up to temperature rescaling.

\emph{Status:} Should be straightforward from the Gaussian KL and Bures formulas.

\item \textbf{Softmax attention as entropy-regularized semi-coupling.}
\emph{Claim:} $\beta_{ij} = \text{softmax}_j(-C_{ij}/\kappa)$ solves $\min_{\beta \geq 0, \beta\mathbf{1}=\mathbf{1}} \sum_{ij}\beta_{ij}C_{ij} + \kappa\sum_{ij}\beta_{ij}\log\beta_{ij}$, i.e., a one-sided Sinkhorn problem.

\emph{Status:} Standard convex optimization result.

\item \textbf{Gauge transport as constrained Monge map.}
\emph{Claim:} For Gaussians, $\Omega_{ij}^* = \arg\min_{\Omega \in \GL(K)} W_2^2(q_i, (\Omega)_\#q_j)$ recovers the Brenier map $A(\Sigma_j, \Sigma_i)$ composed with mean translation.

\emph{Status:} Requires showing that the GL(K)-parameterized family is rich enough to contain the Brenier map (which it is, since Brenier maps between Gaussians are linear).

\item \textbf{Fisher-Rao natural gradient $=$ Wasserstein gradient at isotropic limit.}
\emph{Claim:} When $\Sigma = \sigma^2 I$, the Fisher-Rao natural gradient $\dot{\mu} = -\Sigma\nabla_\mu F = -\sigma^2\nabla_\mu F$ and the Wasserstein gradient $\dot{\mu} = -\nabla_\mu F$ differ only by a scalar factor (the isotropic precision).

\emph{Status:} Immediate.

\item \textbf{GL(K) invariance of gauged Bures distance.}
\emph{Claim:} $W_{2,\Omega}^2(q_i, q_j) = W_2^2(q_i, (\Omega_{ij})_\#q_j)$ is invariant under simultaneous gauge transformation $q_i \mapsto g_\#q_i$, $q_j \mapsto g_\#q_j$, $\Omega_{ij} \mapsto g\Omega_{ij}g^{-1}$ for all $g \in \GL(K)$ if and only if the Bures metric is replaced by the affine-invariant distance.

\emph{Status:} Requires careful computation. The standard Bures metric breaks GL(K) invariance; the affine-invariant (Fisher-Rao) metric preserves it. This is a key tension point.

\end{enumerate}

\subsection{Key Tension: GL(K) Invariance vs.\ Metric Structure}

This is the central obstacle to naive unification:

\begin{center}
\begin{tabular}{lccc}
\toprule
\textbf{Distance} & \textbf{Metric?} & \textbf{GL(K)-invariant?} & \textbf{OT interpretation?} \\
\midrule
$\KL(p\|q)$ & No (asymmetric) & Yes & Entropic divergence \\
$W_2^2(p,q)$ & Yes & No (only $O(K)$) & Quadratic OT \\
Fisher-Rao $d_{\text{FR}}$ & Yes & Yes & Not standard OT \\
Bures $d_{\text{BW}}$ & Yes & No (only $O(K)$) & Gaussian $W_2$ \\
\bottomrule
\end{tabular}
\end{center}

The resolution likely involves one of:
\begin{enumerate}
    \item[(a)] Restricting the gauge group to $O(K)$ when using Bures/Wasserstein costs (sacrificing GL(K) for metric structure).
    \item[(b)] Using the affine-invariant metric $d_{\text{AI}}^2 = \|\log(\Sigma_1^{-1/2}\Sigma_2\Sigma_1^{-1/2})\|_F^2 + \|\mu_1 - \mu_2\|_{\Sigma}^2$ which is GL(K)-invariant and a proper metric, but does not have a standard OT interpretation.
    \item[(c)] Defining a \emph{gauged Wasserstein distance} that absorbs the gauge freedom into the transport cost, preserving both metric structure and gauge covariance.
\end{enumerate}

Option (c) is the most promising and requires the least sacrifice.


%=============================================================================
\section{Computational Implementation Roadmap}
%=============================================================================

\subsection{Phase 1 Implementation: Bures Attention Module}

Add a Bures-Wasserstein attention option alongside the existing KL attention.

\textbf{Required:}
\begin{itemize}
    \item Implement $W_2^2$ between Gaussians (Eq.~\ref{eq:bures}). Requires matrix square root $(\Sigma_i^{1/2}\Sigma_j\Sigma_i^{1/2})^{1/2}$, computable via eigendecomposition.
    \item Implement gauge-covariant Bures: $W_{2,\Omega}^2(q_i, q_j) = \|\mu_i - \Omega_{ij}\mu_j\|^2 + \tr(\Sigma_i + \Omega_{ij}\Sigma_j\Omega_{ij}^\top - 2(\Sigma_i^{1/2}\Omega_{ij}\Sigma_j\Omega_{ij}^\top\Sigma_i^{1/2})^{1/2})$.
    \item Gradients: differentiate through matrix square root (can use implicit differentiation or Sylvester equation).
    \item A/B test: compare KL-attention vs.\ Bures-attention on WikiText at matched parameters.
\end{itemize}

\subsection{Phase 2 Implementation: Sinkhorn Normalization}

\textbf{Required:}
\begin{itemize}
    \item Implement Sinkhorn iterations on the attention cost matrix (5--20 iterations typically suffice).
    \item Support both uniform and learned marginal constraints.
    \item Compare standard softmax attention vs.\ Sinkhorn attention on convergence and perplexity.
\end{itemize}

\subsection{Phase 3 Implementation: JKO Belief Steps}

\textbf{Required:}
\begin{itemize}
    \item Replace the current belief update $\mu \leftarrow \mu - \eta\Sigma\nabla F$ with a JKO proximal step:
    $\mu^{(t+1)} = \arg\min_\mu \{F[\mu] + \frac{1}{2\tau}W_2^2(\mathcal{N}(\mu,\Sigma), \mathcal{N}(\mu^{(t)},\Sigma))\}$.
    \item For Gaussians with fixed $\Sigma$, this reduces to: $\mu^{(t+1)} = \arg\min_\mu \{F[\mu] + \frac{1}{2\tau}\|\mu - \mu^{(t)}\|^2\}$, which is standard proximal gradient descent. The Wasserstein structure becomes non-trivial when $\Sigma$ also evolves.
    \item Implement the coupled $(\mu, \Sigma)$ JKO update.
\end{itemize}


%=============================================================================
\section{Speculative Directions}
%=============================================================================

\subsection{McCann Interpolation as Gauge Geodesic}

In OT, the \emph{McCann interpolation} between measures $\mu_0$ and $\mu_1$ is:
\[
\mu_t = ((1-t)\text{id} + tT)_\#\mu_0
\]
where $T$ is the Brenier map. For Gaussians, this gives:
\[
\mu_t = \mathcal{N}((1-t)m_0 + tm_1,\; ((1-t)I + tA)S_0((1-t)I + tA)^\top)
\]

In the gauge framework, interpolating between belief $q_i$ and transported belief $\Omega_{ij}q_j$ along the gauge geodesic gives a similar structure. The McCann interpolation may be the ``natural'' geodesic on the fiber bundle when the base is equipped with the Bures geometry.

\subsection{Unbalanced OT and Belief Creation/Destruction}

Standard OT assumes mass conservation. \emph{Unbalanced optimal transport} relaxes this, allowing mass to be created or destroyed at a penalty:
\[
\OT_{\text{unbal}}(\mu,\nu) = \inf_\pi \int c\,d\pi + \lambda_1\KL(\pi_1\|\mu) + \lambda_2\KL(\pi_2\|\nu)
\]

This connects to the Gauge-VFE observation term: observations \emph{inject} information into the system (symmetry breaking), which is analogous to mass creation in unbalanced OT. The self-coupling term $\KL(q\|p)$ acts as a ``mass penalty'' keeping beliefs close to priors.

\subsection{Neural OT and Gauge Learning}

Recent work on neural optimal transport (Korotin et al., 2022) parameterizes transport maps with neural networks. The gauge framework parameterizes transport maps with Lie group elements $\Omega = \exp(\phi)$. A hybrid approach could:
\begin{itemize}
    \item Use gauge transport for the \emph{structured} component (rotation, scaling),
    \item Use a residual neural network for the \emph{unstructured} correction,
    \item Train both via the VFE objective.
\end{itemize}


%=============================================================================
\section{Summary: The Unified Picture}
%=============================================================================

\begin{center}
\begin{tabular}{lll}
\toprule
\textbf{Concept} & \textbf{Gauge-VFE} & \textbf{Optimal Transport} \\
\midrule
Objects & Beliefs $q_i = \mathcal{N}(\mu_i, \Sigma_i)$ & Measures $\mu_i$ \\
Cost & $\KL(q_i \| \Omega_{ij}q_j)$ & $c(x,y) = \|x-y\|^2$ (or general) \\
Coupling & $\beta_{ij}$ (softmax attention) & $\pi_{ij}$ (transport plan) \\
Map & $\Omega_{ij} \in \GL(K)$ (gauge transport) & $T: T_\#\mu = \nu$ (Monge map) \\
Regularization & $\kappa$ (attention temperature) & $\varepsilon$ (entropic reg.) \\
Gradient flow & Fisher-Rao natural gradient & Wasserstein gradient (JKO) \\
Hierarchy & RG flow / meta-agents & Multi-scale OT \\
Symmetry & $\GL(K)$ gauge invariance & Isometry invariance \\
Conservation & Not enforced (softmax rows) & Mass conservation (marginals) \\
Geometry (base) & Pullback of Fisher-Rao & Pullback of Bures-Wasserstein \\
\bottomrule
\end{tabular}
\end{center}

The unification thesis: \textbf{Gauge-VFE attention is a gauge-equivariant, entropically regularized, semi-relaxed optimal transport problem on the statistical fiber bundle, with the gauge connection playing the role of a structured Monge map.}

The research program proceeds from concrete (Phase 1: swap cost functions, benchmark) to speculative (Phase 5: prove RG = multi-scale OT), with each phase producing testable predictions and publishable results.

\end{document}
